\subsection{Size Generalization and Cross-State Transfer}
\label{sec:relwork:generalization}

\subsubsection{Size and Distribution Shift in GNNs}
\label{sec:relwork:generalization:shift}
Message-passing weights are shared across nodes and layers, so a trained model accepts a
graph of any size, a point the inductive line established in principle
\citep{hamilton2017graphsage}. Accepting an input is not generalising to it.
\citet{buffelli2022sizeshift} measure accuracy falling as test graphs grow past the sizes
seen in training, and regularise toward size-invariant representations to recover part of
the loss. A reduction moves size and structure together, so the shift it induces between
training and inference is bounded below by the one they report.

\subsubsection{Train-Test Structural Mismatch}
\label{sec:relwork:generalization:mismatch}
The direct evidence is thin and it points both ways. SCAL trains on coarsened graphs and
evaluates on the originals with usable node-level accuracy under a single set of weights
\citep{huang2021scal}. \citet{generale2022scaling} reach the same conclusion on knowledge
graphs, but only after mapping super-node representations back onto their members, a step no
graph-level target provides. CTS-Bench runs the analogous evaluation on netlists and lands
below $R^2 = 0$ (\ref{sec:relwork:domain:circuits}), the one published measurement on
circuit data \citep{cts_bench2026}. One success on a different task, one success behind a
transfer step, and one failure in the closest domain leave the question open on graph-level
circuit regression. A study that asks it therefore needs a compression whose losslessness is
established in advance (\ref{sec:method:architecture:exact}); without one, a poor result
across structural states cannot be separated from a broken evaluation path.
